\chapter{Spacial Confinement}
\label{chap:digital}

\section*{Chapter description}

This chapter presents a preliminary investigation of spatial confinement effects and the electronic structure of Ge–Si digital alloys. The study focuses on layered systems grown along the [100] direction, considering both Si and Ge substrates. After verifying the consistency of the structural parameters obtained from the relaxation and Poisson's ratio calculations, the electronic band structures are analysed along a Brillouin zone path that accounts for both in-plane and out-of-plane directions.

Particular attention is given to the atomic contributions to the electronic states, investigated through atom-projected band structures. The results reveal interesting features, including direct and indirect band gaps, band folding effects, and differences in the electronic structure depending on the substrate material. The analysis suggests that the Ge–Si digital alloy may exhibit spatially dependent electronic properties, potentially resembling those of a quantum well system, with distinct band-gap characteristics along the direction perpendicular to the layers.

The results demonstrate the potential of Ge–Si digital alloys as systems for further studies of spatial confinement and electronic properties in layered materials. They also identify several important questions for future research, including the influence of layer thickness, the effect of growth along the [111] direction, the spatial distribution of band gaps, and the optical properties associated with possible electronic transitions. The possibility of extending this approach to digital alloys with wurtzite structures is also discussed.

\ornament

\section{Atomic Positions}

In the beginning, Poisson's ratio test results were evaluated.
The difference of atomic positions and lattice vectors between relaxation procedure and Poisson's ratio on testing structures prove to be on the fourth decimal point or below.
This is sufficient enough to use Poisson's ratio and atomic position of previous calculations instead of relaxation procedure.
It should be stated here that this was possible only due to the usage of the same \XC\ functional and pseudopotential  as it was practiced in the investigations of the strain effects (Chpt. \ref{chap:strains}).

\ornament

\section{Band structure}

Although the bulk crystal structures of the component materials are the the diamond ones, due to layered nature, the \BZ\ of the system differs (as it has been explained above).
It is not known to which points the focus should be given so the path was chosen to cover all of the characteristic points and to be the most representative.
For the [100] direction the path was set to $Y-\Gamma-X-U-Z-\Gamma$.
The crucial criterium was to cover both in--plane and along crystal growth directions to analyze how the electrons will behave in both directions.

\begin{figure}[!htbp]
	\centering
	\subfloat[][Band structure.]{\includegraphics[width=.96\textwidth]{/digital/bandSi.pdf}}\\
	\subfloat[][$Ge_1$]{\includegraphics[width=.48\textwidth]{/digital/fatband1Si.pdf}}\quad
	\subfloat[][$Ge_2$]{\includegraphics[width=.48\textwidth]{/digital/fatband2Si.pdf}}\\
	\subfloat[][$Si_1$]{\includegraphics[width=.48\textwidth]{/digital/fatband3Si.pdf}}\quad
	\subfloat[][$Si_2$]{\includegraphics[width=.48\textwidth]{/digital/fatband4Si.pdf}}
	\caption{Atomic decomposition of band structure of [100] direction on Si substrate.}
	\label{Fig:da_fbnd_si_1_100}
\end{figure}

\begin{figure}[!htbp]
	\centering
	\subfloat[][Band structure.]{\includegraphics[width=.96\textwidth]{/digital/bandGe.pdf}}\\
	\subfloat[][$Ge_1$]{\includegraphics[width=.48\textwidth]{/digital/fatband1Ge.pdf}}\quad
	\subfloat[][$Ge_2$]{\includegraphics[width=.48\textwidth]{/digital/fatband2Ge.pdf}}\\
	\subfloat[][$Si_1$]{\includegraphics[width=.48\textwidth]{/digital/fatband3Ge.pdf}}\quad
	\subfloat[][$Si_2$]{\includegraphics[width=.48\textwidth]{/digital/fatband4Ge.pdf}}
	\caption{Atomic decomposition of band structure of [100] direction on Ge substrate.}
	\label{Fig:da_fbnd_ge_1_100}
\end{figure}

As a preliminary results, in Figs. \ref{Fig:da_fbnd_si_1_100} and \ref{Fig:da_fbnd_ge_1_100} the band structures are shown for the cases of (100) plane and Si and Ge substrates, together with the decomposition for each atom in elementary cell. Some interesting features of the band structures can be observed. In both systems the direct gap of the order of $1eV$ can be observed, which is presumably a result of  Si $\Gamma$-$X$ band folding. However, in the case of Si substrate, the fundamental gap seems to be indirect, from $\Gamma$ point to $U$ point, but the situation looks different when looking at the atom projected band structures. Here, considering the Ge2 and Si1 atoms, the band gap is direct. Moreover, the $X-U$ direction is perpendicular to the layered system, and a question arises about spatial distribution of band gap energies in the direction perpendicular to the layered systems. Here, a kind of quantum well system can be expected, which can be studied by evaluating the spatially resolved band gaps, for the vertical direction. In the case of Ge substrate the band gap is direct around $\Gamma$, is smaller than $1eV$, and in the direction $\Gamma-Z$ (perpendicular to the layers) has a nesting character, wheres at $U$ point the bands at CB minimum are degenerate. 

The presented results show that the idea of the Ge-Si digital alloy is promising and deserves further attention.  Some interesting features of electronic band structures have been demonstrated here, leading to important questions which should be addressed: 1. what is the effect of increasing number of a given material layers, 2. how the band structures will look like for (111) direction, 3. how the spatially resolved band gaps (in the  vertical direction) behave, and what are the consequences of this behavior for optical properties, 4. what are the oscillator strengths of identified possible optical transitions, and how they, together with the energetic structure affect the optical response. Another intriguing  perspective is an extension of the  investigations on the Ge-Si digital alloys in wurtzite structures, discussed in Chpt. \ref{chap:wurtzite}, where germanium exhibits already a direct fundamental gap, independently of strains, and the role silicon should be determined. Here, the c-direction seems to be natural to form such an alloy, and the band folding in this direction can introduce unpredicted effects. All those issues will be a subject of the planned further investigations.

\ornament